% ============================================================
% SECTION 4.6: THE BOUNDARY. Written fresh 2026-07-26 (late) from the CSVs.
% This REPLACES sections/05_boundary.tex, which Mohammad retired: it predated
% the protocol section, and three of its numbers did not survive checking.
%
% Numbers verified in this pass, with what changed from the retired section:
%   LFR counts by degree      exp_AA via exp_AC/analyse.py LFR reference block
%   seed robustness           RECOMPUTED from exp_AA/metis_noise.csv. The retired
%                             text said the effect clears noise "by a factor of
%                             between two and sixty". True range at d>=49 is
%                             -1.32 to 5.79 in modularity and -1.88 to 14.36 in
%                             AMI, and several cells at retention 0.5 are
%                             negative. Restated per retention point.
%   "harmful on both measures below 25"  CONTRADICTED by exp_AA's own table
%                             (6/12 positive on AMI at 24.6). Removed.
%   speed                     RECOMPUTED from exp_AA/results_armB.csv: median
%                             pipeline speedup 0.75x at d=10 falling to 0.26x at
%                             d=200. The retired text's 1.49x / 1.01x / 0.58x
%                             does not match arm B.
%   edge-budget control       DSpar is positive in 2/12 and 3/12 cells at the two
%                             highest degrees even on LFR, and on the one real
%                             network above the transition it reproduces the gain
%                             outright (exp_AD). "Does not reproduce at any
%                             degree" is retired.
%   Leiden arm                exp_AA/results_armA.csv: 28/144 lspar cells beat
%                             best-of-5, none above +0.005; resolution-matched
%                             control wins 104/144.
%   generalization            exp_AC/SUMMARY.md and analysis_full.txt
%   real networks             exp_AD/results.csv
%   denoising                 exp_AA/results_armC.csv, leiden+lspar 4/4 monotone
% STYLE: no em-dashes, no AI-characteristic phrasing (see STYLE.md).
% ============================================================

\subsection{The boundary}\label{sec:experiments:boundary}

Everything so far is measured on one side of a line. The networks in
Table~\ref{tab:networks} have average degrees between $5.5$ and $32.6$, and the
detectors used for most of it choose their own granularity. Both conditions in
this paper's claim are therefore held at one setting, and a negative result
obtained that way says nothing about where the setting stops mattering. This
subsection varies both.

\paragraph{What the setting requires}
Three properties are needed. Density must vary over an order of magnitude while
community structure is held fixed, which real networks cannot provide. Community
membership must be known rather than inferred. And the granularity artifact of
Section~\ref{sec:experiments:accounting} must be removed by construction rather
than by control, which a detector taking $k$ as an input does: both arms return
the same number of communities by definition.

\paragraph{The transition}
Table~\ref{tab:boundary} reports, for LFR graphs at each average degree, the
number of configurations in which L-Spar improves modularity measured on the
original graph under Metis, and the number in which it improves chance-corrected
agreement with the planted labels.

\begin{table}[t]
\centering
\caption{Metis with $k$ pinned to the planted community count, LFR, realized
mixing near $0.6$. Counts are configurations out of twelve in which L-Spar
improves the quantity named, with the mean change beside them. Modularity is
measured on the original graph. Because $k$ is an input, both arms return
identical community counts and no granularity effect is possible.}
\label{tab:boundary}
\small
\setlength{\tabcolsep}{5pt}
\begin{tabular}{@{}lcccc@{}}
\toprule
$d_{\mathrm{avg}}$ & $\Delta Q > 0$ & mean $\Delta Q$ &
$\Delta\mathrm{AMI} > 0$ & mean $\Delta\mathrm{AMI}$\\
\midrule
11.2  & 0/12  & $-0.103$ & 0/12  & $-0.216$\\
24.6  & 0/12  & $-0.046$ & 6/12  & $-0.072$\\
49.4  & 8/12  & $-0.016$ & 9/12  & $+0.042$\\
101.9 & 10/12 & $+0.006$ & 12/12 & $+0.116$\\
220.9 & 12/12 & $+0.009$ & 12/12 & $+0.093$\\
\bottomrule
\end{tabular}
\end{table}

The two measures cross at different places. Recovery turns positive first, with
half the configurations improving at average degree $24.6$ and three quarters at
$49.4$, while the objective is still negative on average there and only becomes
positive in the mean at $101.9$. The location of the crossing was not chosen by
us and was not known to us when the experiment was designed: it coincides with
the degree at which the original authors reported their method beginning to
overtake the unsparsified baseline.

The size of the effect should be read with the counts. At the degree where the
objective crosses, the mean change is still negative and the positive cells are
worth thousandths of a modularity point. The recovery gains are an order of
magnitude larger, reaching a mean of $+0.116$.

\paragraph{Run-to-run variation}
Metis is not deterministic across option seeds, so the decisive cells were
repeated with ten seeds per graph and the worst sparsified run compared against
the best unsparsified run. The comparison separates the two measures again. On
recovery it is comfortable: at aggressive retention the worst case clears the
baseline seed spread by factors of $4.9$ to $19.8$, and it clears it at every
degree at or above $49$. On the objective it is not: the same cells reach factors
of only $7.8$ down to $-0.4$, so one of them fails to clear at all, and at the
mildest retention the cells at $201$ and $261$ are negative in the worst case on
both measures. The recovery result survives a hostile reading of the seeds; the
modularity result survives at the aggressive operating points with one exception.

\paragraph{Generality across generators}
LFR graphs are one family with one degree distribution, and the mechanism runs
through a degree-derived or similarity-derived score, so the family could be
doing the work. Repeating the arm on stochastic block models and on
degree-corrected block models, calibrated cell by cell to the LFR graphs'
realized degree and mixing, reproduces the sign flip on both. It also refutes the
explanation we expected. We predicted that a degree-homogeneous generator would
push the transition to higher density, since a degree-based selector has less to
work with, and that the degree-corrected variant would sit closer to LFR. Neither
holds. All four block-model variants cross at essentially the same place while
their degree coefficient of variation spans an order of magnitude: at the degree
where the transition sits it runs from $0.14$ to $1.47$ across the four
generators. Across sixty generator-by-degree cells the gain correlates with
average degree at Spearman $+0.65$ and with degree variability at $-0.26$. Degree
heterogeneity is not the operative variable.

\paragraph{Real networks}
The same arm on the seven real networks places the transition without resolving
what sits on the far side of it. Below average degree $11$ there is nothing: 128
cells across five networks, seven sparsifiers, and both ways of choosing $k$ on
the two of those networks that carry reference communities, with not
one positive in the worst case over ten seeds. On wiki-Vote, at average degree
$28.5$, four of the seven sparsifiers produce a worst-case positive gain, the
largest being $+0.0124$. On email-Eu-core, at $32.6$, none does: two arms are
positive in the mean and every worst case is negative. The transition on real
graphs therefore lies between $10.7$ and $28.5$, and above it the effect is
present on one of the two networks we can test.

\paragraph{What produces the gain}
On the benchmark graphs the selection rule matters: degree-based and uniform
sparsification at identical retention are negative at every degree except the two
highest, where DSpar turns positive in 2 and 3 of twelve cells against L-Spar's
10 and 12. On wiki-Vote that separation disappears. The four arms that score
individual edges locally all produce the gain, DSpar among them at $+0.0076$
against L-Spar's $+0.0124$, and the two connectivity-preserving backbones destroy
it. What the gain requires is a rule that scores edges one at a time; which
signal the rule uses matters less than we expected, and we do not identify the
property that separates Local Degree, which is also local and is the worst arm on
that network.

\paragraph{Why density}
The denoising account makes a testable prediction. If sparsification helps by
deleting edges that do not belong to the generative structure, its benefit should
grow with the proportion of such edges. Injecting uniformly random edges into
benchmark graphs while holding the planted labels fixed is consistent with it for
similarity-based selection: recovery gain rises with the injected
fraction in three of the four configurations tested and rises with one reversal
in the fourth, and at average degree
$50$ with $20\%$ retention it grows from $+0.061$ with no injected noise to
$+0.185$ when the number of spurious edges equals the number of real ones. The
gain does not vanish at zero noise, so denoising adds to an effect already
present rather than being its sole cause, and a benchmark's noise is random by
construction whereas a measurement process need not be. The experiment
establishes the direction, not the magnitude in the field.

\paragraph{The same graphs, the other detector family}
Running a free-granularity optimizer on the identical graphs gives a different
answer, which is what makes the boundary a boundary rather than a threshold in
density alone. Honest modularity gains appear in 28 of 144 configurations and
none exceeds $+0.005$, within the spread of five random restarts. Recovery gains
do appear and grow with degree, but they do not survive the granularity control:
partitioning the original graph at a resolution tuned to the sparsified
partition's community count recovers the planted communities as well or better in
104 of the 144 configurations where that comparison is defined. What looked like
a benefit of sparsification is the benefit of a finer partition, available
without touching the graph.

\paragraph{Three things the positive result does not say}
Sparsification repairs a weaker detector rather than producing the best available
partition. At average degree $49$, Metis with L-Spar reaches AMI $0.75$ against
its own baseline of $0.62$, a real improvement to that algorithm, while a
resolution-tuned optimizer on the untouched graph reaches $0.94$ on the same
graphs. The defensible statement is that similarity-based selection repairs a
fixed-$k$ cut partitioner on dense graphs, not that it improves community
detection.

The quality gain is not a speed gain in our implementation. Computing exact
pairwise similarity costs more than it saves, and the deficit grows with density
because the similarity computation grows faster than the detection it
accelerates: the median end-to-end speedup falls from $0.89\times$ at average
degree $11$ to $0.18\times$ at $221$. The original speedups were obtained with an
approximate similarity computation reported as two orders of magnitude cheaper,
measured against clustering runs lasting hours, so both cost bases differ from
ours by enough to explain the discrepancy without disagreement about what should
be counted.

One part of the original report does not reproduce. At high mixing, where the
planted structure is barely detectable and baseline recovery sits between $0.03$
and $0.17$, no gain appears at any degree: across 104 configurations L-Spar
improves modularity in 9, and the mean change is negative at every degree. The
original sweep reports its largest improvement in that regime. Whether this is a
genuine non-reproduction or a floor effect of the metrics we use cannot be
settled at that mixing level, and we report it unresolved.
